\scanpage{515}
\mspage{426}
% B01
\uline{Corollaire} 4.15.2. Les hypothèses sur $f:E\unsure{\relbar|\longrightarrow}C$ sont celles de 4.15.1 ($C$ petite, et $f$ satisfait $H'=H1)2)3)$).

% B02
% slanted
\margin{\inserted{Pour a) et b),} On n'utilise que $H1$ et $H2$, à l'exclusion de $H3$.}

% B03
a) Soit $s\in\mathrm{Ob}\,C$ et soit $\pi_1$ un cardinal infini tel que pour tt $\lambda:s\to t$ (flèche de $C$ de source $t$), on ait $\pi_1\geq\pi_\lambda$ (où $\pi_\lambda$ est défini dans l'énoncé de $H2$). Alors le foncteur $i_s:E_s\longrightarrow E$ \deleted{satisfait} est $\pi_1$-accessible.

% B04
b) \inserted{On suppose $C$ \uline{sans projecteurs}.} Soit $\pi_{\deleted{\unsure{0}}1}$ un cardinal \deleted{\unsure{infini tel que,} 1, les catégories $E_s$ ($s\in\mathrm{Ob}\,C$) satisfassent $L_{\pi_1}$ (existe car : $H1)$ et au fait que $C$ ou $\mathrm{Ob}\,C$ petit)), et tel que $\pi_1\geq\mathrm{card}\,\deleted{\unsure{\mathrm{Ob}}}C$.} \inserted{$\geq\mathrm{Sup}(\underset{\lambda}{\mathrm{Sup}}\pi_\lambda,\mathrm{card}\,C)$} Alors $E$ satisfait $L_{\pi_1}$, et $f:E\to C$ est $\pi$-accessible.

% B05
c) \deleted{Soit $\pi_0$ un \unsure{cardinal}} \inserted{Supposons} que $C$ soit sans projecteurs. (p. ex une catégorie ``rigide'', p. ex. une catégorie \inserted{pré}ordonnée). Soient $\pi_0,\pi_1$ des cardinaux tels que

% B06
\textcircled{\scriptsize $1^\circ$} \deleted{$\pi_0\geq\unsure{\sup}$ \unsure{que}} $\forall s\in\mathrm{Ob}\,C$, $\pi_0$ est \uline{utile pour} $E_s$ (\inserted{un tel $\pi_0$} existe grâce à $H1)$ et au fait que $C$ \deleted{\unsure{est un}} $\mathrm{Ob}\,C$ petit).

% B07
\textcircled{\scriptsize $2^\circ$} $\pi_1\geq\mathrm{Sup}(\pi_0,\inserted{\mathrm{card}\,C,}\underset{\lambda\in\mathrm{Fl}\,C}{\mathrm{Sup}}\pi_\lambda,\underset{s\in\mathrm{Ob}\,C}{\mathrm{Sup}}\underline{\mathrm{card\,red}}\ E_{s\pi_0})$.

% B08
\textcircled{\scriptsize $3^\circ$} $\forall\lambda\in\mathrm{Fl}\,C$, et $x_0\in\mathrm{Ob}\,E_{s,\pi_0}$, le foncteur $\unsure{H_\lambda|h_\lambda|h} (x_0)$ de $H3$ \inserted{$x_1\longmapsto\mathrm{Hom}_{\lambda}(x_0,x_1): E_t\longrightarrow\mathrm{Ens}$ est $\pi_1$-\uline{accessible}.}
% slanted

% B09
Soit $\pi$ \deleted{\unsure{$\geq\pi_1$}} un cardinal tel que $\pi^{\pi_0}=\pi\geq\pi_1$. Alors Par b) $E$ satisfait $L_{\pi_1}$ donc $L_\pi$, et $f:E\to C$ est $\pi_1$-accessible\deleted{\unsure{ment}}, De plus
% B10
\[i_s(E_{s\pi})\subset E_\pi,\]
% B11
\deleted{d'où} ou encore \deleted{et \unsure{$E_\pi$} d'où un foncteur pl. fid ($E_\pi$ satisfaisant à $L_{\pi_1}$, et a fortiori à $L_\pi$, grâce à b)} \inserted{(cf)}
% B12
\[E(\pi)\subset E_\pi,\]
% B13
où $E(\pi)$ est la sous-catégorie pl. de $E$ \deleted{formée des} telle que
% B14
\[\mathrm{Ob}\,E(\pi)=\bigcup_{s\in\mathrm{Ob}\,C}\mathrm{Ob}(E_{s\pi}).\]

% B15
Comme $E_\pi$ satisfait $L_{\pi_1}$ par b), donc a fortiori $L_\pi$, on a donc en ce cas un foncteur \uline{pl fid.}
% B16
\[\mathrm{Ind}_{\pi}(E_{s\pi})\xrightarrow{\approx}E_{\deleted{\unsure{s|\pi}}},\]
% B17
et ce foncteur est une équivalence de catégories (donc $E$ est accessible).

% B18
d) En fait, sous les hyp. de c) sur $C,\pi_0,\unsure{\pi_\lambda|\pi_1},\pi$, on a $E(\pi)=E_\pi$ ! (\unsure{ou|non} seulement $E(\pi)\subset E_\pi$)

% archive sheet number: 514
